\documentclass[12pt,a4paper]{article}
\usepackage[margin=1in]{geometry}
\usepackage{booktabs}
\usepackage{longtable}
\usepackage{array}
\usepackage{hyperref}
\usepackage{titlesec}
\usepackage{parskip}

\hypersetup{
    colorlinks=true,
    linkcolor=blue,
    urlcolor=blue
}

\title{Dataset Card}
\date{}

\begin{document}

\maketitle

% ─────────────────────────────────────────────
\section*{Summary Information}
\textit{This describes the data set that you are going to use and what it was used for.
You may choose to list models in here that used the data described in the card.
The date is also important due to how data may change over time.
Dataset cards do not necessarily describe all of the data provided by a sponsor or for the project,
but they should describe all data \textbf{used} for the project.}

\begin{tabular}{@{}ll@{}}
\toprule
\textbf{Field}       & \textbf{Description} \\
\midrule
Name                 &                      \\
Curation Date        &                      \\
Sensitivity Level    &                      \\
Summary              &                      \\
Source               &                      \\
\bottomrule
\end{tabular}

% ─────────────────────────────────────────────
\section*{Card Authors}
\textit{This describes the people that contributed to creating the dataset card.}

\begin{tabular}{@{}ll@{}}
\toprule
\textbf{Name} & \textbf{Email} \\
\midrule
Sam Freund          & freunds@msoe.edu  \\
\bottomrule
\end{tabular}

% ─────────────────────────────────────────────
\section*{Data Overview}
\textit{General description of the data set. Such as number of features, how large the files are,
number of observations, etc.}

\begin{tabular}{@{}ll@{}}
\toprule
\textbf{Field}      & \textbf{Value} \\
\midrule
Storage Size        &                \\
Number of Rows      &                \\
Number of Features  &                \\
Data Format         &                \\
\bottomrule
\end{tabular}

% ─────────────────────────────────────────────
\section*{Numerical Features Summary}
\textit{This section shows the summary statistics for all numerical features in the data set.}

\begin{longtable}{@{}lllllllll@{}}
\toprule
\textbf{Feature} & \textbf{Count} & \textbf{Mean} & \textbf{Std Dev} &
\textbf{Min} & \textbf{25\%} & \textbf{Median} & \textbf{75\%} & \textbf{Max} \\
\midrule
\endfirsthead
\toprule
\textbf{Feature} & \textbf{Count} & \textbf{Mean} & \textbf{Std Dev} &
\textbf{Min} & \textbf{25\%} & \textbf{Median} & \textbf{75\%} & \textbf{Max} \\
\midrule
\endhead
 & & & & & & & & \\
\bottomrule
\end{longtable}

% ─────────────────────────────────────────────
\section*{Categorical Features Summary}
\textit{This section shows the equivalent summary statistic for all categorical features in the data set.}

\begin{tabular}{@{}lll@{}}
\toprule
\textbf{Feature} & \textbf{Unique Values} & \textbf{Most Common Value} \\
\midrule
                 &                        &                            \\
\bottomrule
\end{tabular}

% ─────────────────────────────────────────────
\section*{Field Information}
\textit{Document the table that is encoded, i.e.\ the table they are using to build the encoding.}

\begin{longtable}{@{}llll@{}}
\toprule
\textbf{Feature Name} & \textbf{Data Type} & \textbf{Statistical Type} & \textbf{Description} \\
\midrule
\endfirsthead
\toprule
\textbf{Feature Name} & \textbf{Data Type} & \textbf{Statistical Type} & \textbf{Description} \\
\midrule
\endhead
Feature\_example & int & Nominal & An example feature created for the use of the dataset card format. \\
\bottomrule
\end{longtable}

% ─────────────────────────────────────────────
\section*{Example Entry}
\textit{It is helpful to track a few observations that you are familiar with.
Pick an ``average'' observation as well as a few that exhibit features or classes that you are
interested in. E.g.\ for a cancer dataset you might choose a sample that has all of the classes
that you want to predict or that shows severe progression of the disease.}

% ─────────────────────────────────────────────
\section*{Exploratory Charts}
\textit{Include a scatter plot matrix of all of your numeric features.
Show bar graphs for your categorical features.
If you know what response you are interested in, you might choose to make a plot of the first two
principal components (from PCA or similar feature reduction algorithm) vs.\ your known response
or classes.}

% ─────────────────────────────────────────────
\section*{Notable Feature Processing}
\textit{If you decide to do feature scaling or if you are doing any feature engineering, the process
should go here. Remember that the mean and standard deviation of features would be important for
any future applications that might use this data.}

% ─────────────────────────────────────────────
\section*{Notes}
\textit{Document any additional considerations, concerns, or otherwise here.}

\end{document}
